\begin{textsample}{NEG Dim 7 – global\_south\_2025\_03 – Score -1.00 – global\_south\_2025\_03/122779006.txt}  \label{ex:f7_neg_027}
CAIRO : Egypt has floated a new proposal aimed at restoring the Gaza ceasefire deal, security sources said on Monday, as Palestinian health authorities said Israeli strikes had killed at least 65 people in the enclave in the past 24 hours.
The proposal, made last week, follows an escalation in violence after Israel resumed air and ground operations against Hamas on March 18, effectively ending a \textbf{two-month} period of relative calm after 15 months of war.
Gaza health officials said Israeli airstrikes and shelling have killed nearly 700 Palestinians since then, including at least 400 women and children.
Among those killed on Monday were two local journalists, Moha? mmad Mansour and Hussam Shabat, medics said.
Hamas said several of its senior political and security officials had also been killed.
The Egyptian plan calls for Hamas to release five Israeli prisoners each week, with Israel implementing the second phase of the ceasefire after the first week, two security sources said.
Both the US and Hamas have agreed @ @ @ @ @ @ @ @ @ @ Israel has not yet responded.
The sources said the Egyptian proposal includes a timeline for a full Israeli military withdrawal from Gaza, backed by US guarantees, in exchange for the release of remaining prisoners.
Education at risk Since Israel launched a major operation in the West Bank in January, the trip has become even more perilous.
Everyday, children in the West Bank run the gauntlet of Israeli roadblocks, checkpoints and settler attacks on their way to school.
Thousands of troops are sweeping through refugee camps and cities and demolishing houses and infrastructure, including roads children use to get to school.
The impact on children ' s education is reminiscent of the havoc caused in the Gaza Strip during the war."The ability for Palestinian children to access quality education in the West Bank or in Gaza has never been under more stress,"said Alexandra Saieh, global head of humanitarian policy and advocacy at Save the Children.
To make matters worse, the UN Palestinian relief agency UNRWA @ @ @ @ @ @ @ @ @ @ East Jerusalem, could be forced to stop its work following an Israeli ban on its operations on Israeli territory.
And funding cuts from major donors, including the United States where President Donald Trump has terminated thousands of foreign aid projects, could further cripple services.
In December, Sweden also cut its support, a decision the agency said came at the worst time for Palestinian refugees.

% matched lemmas: two-month
\end{textsample}
